\documentclass[10pt,a4paper]{amsart}
\usepackage{amsmath,amssymb,amsthm,graphicx}
\usepackage{verbatim}
\newenvironment{reference}{\comment}{\endcomment}
%\newenvironment{reference}{}{}
\newenvironment{slogan}{\comment}{\endcomment}
\newenvironment{history}{\comment}{\endcomment}

% For commutative diagrams we use Xy-pic
\usepackage[all]{xy}

% We use 2cell for 2-commutative diagrams.
\xyoption{2cell}
% Theorem environments.
%
\theoremstyle{plain}
\newtheorem{theorem}[subsection]{Theorem}
\newtheorem{proposition}[subsection]{Proposition}
\newtheorem{lemma}[subsection]{Lemma}

\theoremstyle{definition}
\newtheorem{definition}[subsection]{Definition}
\newtheorem{example}[subsection]{Example}
\newtheorem{exercise}[subsection]{Exercise}
\newtheorem{situation}[subsection]{Situation}

\theoremstyle{remark}
\newtheorem{remark}[subsection]{Remark}
\newtheorem{remarks}[subsection]{Remarks}

\numberwithin{equation}{subsection}

% Macros
%
\def\lim{\mathop{\mathrm{lim}}\nolimits}
\def\colim{\mathop{\mathrm{colim}}\nolimits}
\def\Spec{\mathop{\mathrm{Spec}}}
\def\Hom{\mathop{\mathrm{Hom}}\nolimits}
\def\Ext{\mathop{\mathrm{Ext}}\nolimits}
\def\SheafHom{\mathop{\mathcal{H}\!\mathit{om}}\nolimits}
\def\SheafExt{\mathop{\mathcal{E}\!\mathit{xt}}\nolimits}
\def\Sch{\mathit{Sch}}
\def\Mor{\mathop{\mathrm{Mor}}\nolimits}
\def\Ob{\mathop{\mathrm{Ob}}\nolimits}
\def\Sh{\mathop{\mathit{Sh}}\nolimits}
\def\NL{\mathop{N\!L}\nolimits}
\def\CH{\mathop{\mathrm{CH}}\nolimits}
\def\proetale{{pro\text{-}\acute{e}tale}}
\def\etale{{\acute{e}tale}}
\def\QCoh{\mathit{QCoh}}
\def\Ker{\mathop{\mathrm{Ker}}}
\def\Im{\mathop{\mathrm{Im}}}
\def\Coker{\mathop{\mathrm{Coker}}}
\def\Coim{\mathop{\mathrm{Coim}}}

% Boxtimes
%
\DeclareMathSymbol{\boxtimes}{\mathbin}{AMSa}{"02}

%
% Macros for moduli stacks/spaces
%
\def\QCohstack{\mathcal{QC}\!\mathit{oh}}
\def\Cohstack{\mathcal{C}\!\mathit{oh}}
\def\Spacesstack{\mathcal{S}\!\mathit{paces}}
\def\Quotfunctor{\mathrm{Quot}}
\def\Hilbfunctor{\mathrm{Hilb}}
\def\Curvesstack{\mathcal{C}\!\mathit{urves}}
\def\Polarizedstack{\mathcal{P}\!\mathit{olarized}}
\def\Complexesstack{\mathcal{C}\!\mathit{omplexes}}
% \Pic is the operator that assigns to X its picard group, usage \Pic(X)
% \Picardstack_{X/B} denotes the Picard stack of X over B
% \Picardfunctor_{X/B} denotes the Picard functor of X over B
\def\Pic{\mathop{\mathrm{Pic}}\nolimits}
\def\Picardstack{\mathcal{P}\!\mathit{ic}}
\def\Picardfunctor{\mathrm{Pic}}
\def\Deformationcategory{\mathcal{D}\!\mathit{ef}}
\usepackage{cleveref}
\crefname{originaldefinition}{original definition}{original definitions}
\Crefname{originaldefinition}{Original definition}{Original definitions}
\makeatletter
\newlabel{categories-definition-directed}{{002V}{source}{}{sourcepointer-1}{}}
\newlabel{categories-definition-directed@cref}{{[originaldefinition][0][]002V}{[1][0][]source}{}{}{}}
\makeatother
\crefname{originalsection}{original section}{original sections}
\Crefname{originalsection}{Original section}{Original sections}
\makeatletter
\newlabel{categories-section-directed-colimits}{{04AX}{source}{}{sourcepointer-2}{}}
\newlabel{categories-section-directed-colimits@cref}{{[originalsection][0][]04AX}{[1][0][]source}{}{}{}}
\makeatother
\crefname{originallemma}{original lemma}{original lemmas}
\Crefname{originallemma}{Original lemma}{Original lemmas}
\makeatletter
\newlabel{lemma-independent-refinement}{{00W7}{source}{}{sourcepointer-3}{}}
\newlabel{lemma-independent-refinement@cref}{{[originallemma][0][]00W7}{[1][0][]source}{}{}{}}
\makeatother
\crefname{originallemma}{original lemma}{original lemmas}
\Crefname{originallemma}{Original lemma}{Original lemmas}
\makeatletter
\newlabel{lemma-common-refinement}{{00W6}{source}{}{sourcepointer-4}{}}
\newlabel{lemma-common-refinement@cref}{{[originallemma][0][]00W6}{[1][0][]source}{}{}{}}
\makeatother
\title{Properties of the plus construction and of its iteration}
\author{The Stacks Project Authors}
\date{stacks-project:sites.tex:local-snapshot-20261012; GNU FDL 1.2 or later}
\begin{document}
\maketitle
\noindent\textit{Editorial note.} Revision note: the first assembly attempt (r001) failed because the counted context span ended one line before the \textbackslash{}end\{definition\} of Definition definition-separated; the span now runs to source line 2189 so that the environment is balanced, and the failed r001 job is preserved as evidence.\par\medskip
\noindent\textit{Allowed standard prerequisites (editorial).}\begin{itemize}
\item The plus construction and its surjectivity. Same chapter, the construction of the plus construction as a filtered colimit and Lemma lemma-plus-surjective. Supplies the description of the elements and the local surjectivity used in the proofs.
\item The definition of a separated presheaf. Same chapter, Definition definition-separated. Supplies the notion appearing in the second statement.
\item Common refinements of coverings. Same chapter, Lemma lemma-common-refinement. Supplies the common refinement used in the first part.
\end{itemize}\medskip
\ifdefined\hypertarget\hypertarget{sourcepointer-1}{}\fi
\noindent\textit{Source reference.} Original definition 002V: sources/stacks-sites/pinned-original-source.tex, source lines 2147--2147. Navigation pointer invoked at proof line 2147; the counted span contains the pointer only and the cited result is not reproduced inside it, so the pointer is registered as a non-load-bearing navigation stub in the established form for Stacks material. No bibliography citation is involved.\par\medskip
\ifdefined\hypertarget\hypertarget{sourcepointer-2}{}\fi
\noindent\textit{Source reference.} Original section 04AX: sources/stacks-sites/pinned-original-source.tex, source lines 2149--2149. Navigation pointer invoked at proof line 2149; the counted span contains the pointer only and the cited result is not reproduced inside it, so the pointer is registered as a non-load-bearing navigation stub in the established form for Stacks material. No bibliography citation is involved.\par\medskip
\ifdefined\hypertarget\hypertarget{sourcepointer-3}{}\fi
\noindent\textit{Source reference.} Original lemma 00W7: sources/stacks-sites/pinned-original-source.tex, source lines 2175--2175. Navigation pointer invoked at proof line 2175; the counted span contains the pointer only and the cited result is not reproduced inside it, so the pointer is registered as a non-load-bearing navigation stub in the established form for Stacks material. No bibliography citation is involved.\par\medskip
\ifdefined\hypertarget\hypertarget{sourcepointer-4}{}\fi
\noindent\textit{Source reference.} Original lemma 00W6: sources/stacks-sites/pinned-original-source.tex, source lines 2220--2220. Navigation pointer invoked at proof line 2220; the counted span contains the pointer only and the cited result is not reproduced inside it, so the pointer is registered as a non-load-bearing navigation stub in the established form for Stacks material. No bibliography citation is involved.\par\medskip
\noindent
From these two lemmas, and the fact that $\mathcal{J}_U$ is nonempty,
it follows that the diagram $H^0(-, \mathcal{F}) : \mathcal{J}_U^{opp}
\to \textit{Sets}$ is filtered, see
Categories, Definition \ref{categories-definition-directed}.
Hence, by Categories,
Section \ref{categories-section-directed-colimits}
the colimit $\mathcal{F}^{+}(U)$ may be described
in the following straightforward manner. Namely, every element in the set
$\mathcal{F}^{+}(U)$ arises from an element
$s \in H^0(\mathcal{U}, \mathcal{F})$ for some covering
$\mathcal{U}$ of $U$. Given a second element $s' \in
H^0(\mathcal{U}', \mathcal{F})$ then $s$ and $s'$ determine
the same element of the colimit if and only if there exists a covering
$\mathcal{V}$ of $U$ and refinements $f : \mathcal{V} \to \mathcal{U}$ and
$f' : \mathcal{V} \to \mathcal{U}'$ such that $f^*s = (f')^*s'$
in $H^0(\mathcal{V}, \mathcal{F})$. Since the trivial covering
$\{\text{id}_U\}$ is an object of $\mathcal{J}_U$ we get
a canonical map $\mathcal{F}(U) \to \mathcal{F}^+(U)$.

\begin{lemma}
\label{lemma-plus-surjective}
The map $\theta : \mathcal{F} \to \mathcal{F}^+$ has the following
property: For every object $U$ of $\mathcal{C}$ and every section
$s \in \mathcal{F}^+(U)$ there exists a covering $\{U_i \to U\}$
such that $s|_{U_i}$ is in the image of $\theta : \mathcal{F}(U_i)
\to \mathcal{F}^{+}(U_i)$.
\end{lemma}

\begin{proof}
Namely, let $\{U_i \to U\}$ be a covering such that $s$ arises
from the element $(s_i) \in H^0(\{U_i \to U\}, \mathcal{F})$.
According to Lemma \ref{lemma-independent-refinement} we may
consider the covering $\{U_i \to U_i\}$ and the (obvious) morphism
of coverings $\{U_i \to U_i\} \to \{U_i \to U\}$ to compute the
pullback of $s$ to an element of $\mathcal{F}^+(U_i)$. And indeed,
using this covering we get exactly $\theta(s_i)$ for the restriction
of $s$ to $U_i$.
\end{proof}

\begin{definition}
\label{definition-separated}
We say that a presheaf of sets $\mathcal{F}$ on a site
$\mathcal{C}$ is {\it separated} if, for all coverings
of $\{U_i \rightarrow U\}$, the map
$\mathcal{F}(U) \to \prod \mathcal{F}(U_i)$ is injective.
\end{definition}

\begin{theorem}
\label{theorem-plus}
With $\mathcal{F}$ as above
\begin{enumerate}
\item
\label{item-sep}
The presheaf $\mathcal{F}^+$ is separated.
\item
\label{item-sheaf}
If $\mathcal{F}$ is separated, then $\mathcal{F}^+$ is a sheaf
and the map of presheaves $\mathcal{F} \to \mathcal{F}^+$ is injective.
\item
\label{item-plus-iso}
If $\mathcal{F}$ is a sheaf, then $\mathcal{F} \to \mathcal{F}^+$
is an isomorphism.
\item
\label{item-plusplus}
The presheaf $\mathcal{F}^{++}$ is always a sheaf.
\end{enumerate}
\end{theorem}

\begin{proof}
Proof of (\ref{item-sep}).
Suppose that $s, s' \in \mathcal{F}^+(U)$ and suppose that
there exists some covering $\{U_i \to U\}$ such that
$s|_{U_i} = s'|_{U_i}$ for all $i$. We now have three coverings
of $U$: the covering $\{U_i \to U\}$ above, a covering $\mathcal{U}$
for $s$ as in Lemma \ref{lemma-plus-surjective},
and a similar covering $\mathcal{U}'$ for $s'$. By Lemma
\ref{lemma-common-refinement}, we can find a common refinement,
say $\{W_j \to U\}$. This means we have $s_j, s'_j \in \mathcal{F}(W_j)$
such that $s|_{W_j} = \theta(s_j)$, similarly for $s'|_{W_j}$, and
such that $\theta(s_j) = \theta(s'_j)$. This last equality means
that there exists some covering $\{W_{jk} \to W_j\}$ such that
$s_j|_{W_{jk}} = s'_j|_{W_{jk}}$. Then since $\{W_{jk} \to U\}$
is a covering we see that $s, s'$ map to the same element of
$H^0(\{W_{jk} \to U\}, \mathcal{F})$ as desired.

\medskip\noindent
Proof of (\ref{item-sheaf}). It is clear that $\mathcal{F} \to
\mathcal{F}^+$ is injective because all the maps
$\mathcal{F}(U) \to H^0(\mathcal{U}, \mathcal{F})$
are injective. It is also clear that, if $\mathcal{U} \to
\mathcal{U}'$ is a refinement, then $H^0(\mathcal{U}', \mathcal{F})
\to H^0(\mathcal{U}, \mathcal{F})$ is injective. Now,
suppose that $\{U_i \to U\}$ is a covering, and let
$(s_i)$ be a family of elements of $\mathcal{F}^+(U_i)$
satisfying the sheaf condition
$s_i|_{U_i \times_U U_{i'}} = s_{i'}|_{U_i \times_U U_{i'}}$
for all $i, i' \in I$. Choose coverings (as in
Lemma \ref{lemma-plus-surjective}) $\{U_{ij} \to U_i\}$
such that $s_i|_{U_{ij}}$ is the image of the (unique)
element $s_{ij} \in \mathcal{F}(U_{ij})$. The sheaf condition
implies that $s_{ij}$ and $s_{i'j'}$ agree over
$U_{ij} \times_U U_{i'j'}$ because it maps to
$U_i \times_U U_{i'}$ and we have the equality there.
Hence $(s_{ij}) \in H^0(\{U_{ij} \to U\}, \mathcal{F})$
gives rise to an element $s \in \mathcal{F}^+(U)$. We leave
it to the reader to verify that $s|_{U_i} = s_i$.

\medskip\noindent
Proof of (\ref{item-plus-iso}). This is immediate from the definitions
because the sheaf property says exactly that every map
$\mathcal{F} \to H^0(\mathcal{U}, \mathcal{F})$ is bijective
(for every covering $\mathcal{U}$ of $U$).

\medskip\noindent
Statement (\ref{item-plusplus}) is now obvious.
\end{proof}

\end{document}
